\documentclass[11pt,a4paper]{article}
\usepackage{ftm}

\title{P versus NP\\[4pt]
       \large Caderno de Lacunas Formais — Programa FTM}
\author{Federated Trivalent Manifold Research Programme}
\date{2026}

\begin{document}
\maketitle
\tableofcontents
\bigskip

%% ---------------------------------------------------------------
\section{Enunciado oficial}
%% ---------------------------------------------------------------

\begin{definition}[P versus NP]
Seja $P$ a classe de linguagens decidíveis em tempo polinomial
por máquina de Turing determinística, e $NP$ a classe com
verificação polinomial. O problema pergunta se
\[
  P = NP.
\]
Uma linguagem $L\in NP$ admite um verificador $V$ em tempo
polinomial: $x\in L \Leftrightarrow \exists w\,V(x,w)=1$
com $|w|=\mathrm{poly}(|x|)$.
\end{definition}

\bigskip

%% ---------------------------------------------------------------
\section{Programas de pesquisa e encaixe FTM}
%% ---------------------------------------------------------------

\begin{programs}

\item \textbf{Barreiras conhecidas.}
Três barreiras demonstram que técnicas usuais não resolvem P vs NP:
\begin{enumerate}[label=(\roman*),leftmargin=2em]
  \item \emph{Relativização} (Baker–Gill–Solovay 1975):
    existem oráculos $A,B$ com $P^A=NP^A$ e $P^B\neq NP^B$.
  \item \emph{Natural Proofs} (Razborov–Rudich 1994):
    qualquer prova de limite inferior via propriedades ``naturais''
    colide com existência de pseudoaleatório.
  \item \emph{Algebrização} (Aaronson–Wigderson 2009):
    técnicas que algebrizam não separam $P$ de $NP$.
\end{enumerate}
Uma separação deve ultrapassar as três barreiras simultaneamente.

\begin{gap}
Não existe framework que supere as três barreiras ao mesmo tempo.
Qualquer abordagem FTM deve identificar qual barreira transgride
e por que a transgressão é legítima. \OPEN
\end{gap}

\item \textbf{Geometric Complexity Theory (GCT).}
Mulmuley–Sohoni propõem separar $P$ de $NP$ via teoria de
representações e obstruções:
\[
  \mathrm{perm}(X)\notin\overline{GL_n\cdot\mathrm{det}(Y)},
\]
onde $\mathrm{perm}$ é o permanente e $\mathrm{det}$ o determinante.
O análogo FTM são obstruções espectrais ao decaimento de
coeficientes de representação — não $W_4$ diretamente.

\begin{gap}[Obstruções GCT]
A existência de obstruções de multiplicidade suficientes para
implementar a separação geométrica permanece conjectura.
O programa GCT produz candidatos, não provas. \OPEN
\end{gap}

\item \textbf{Limites inferiores de circuitos.}
A implicação
\[
  NP\not\subset P/\mathrm{poly}
  \;\Longrightarrow\;
  P\neq NP
\]
reduz o problema a limite inferior em tamanho de circuitos.
Resultados conhecidos:
\begin{itemize}[leftmargin=2em]
  \item $\mathrm{PARITY}\notin AC^0$ (Furst–Saxe–Sipser; Håstad);
  \item $NEXP\not\subset ACC^0$ (Williams 2011).
\end{itemize}
$NP\not\subset P/\mathrm{poly}$ permanece aberto para circuitos gerais.

\begin{gap}
Não há limite inferior superpolinomial provado para circuitos
booleanos sem restrições. \OPEN
\end{gap}

\item \textbf{Derandomização.}
A questão $BPP=P$ (conjectura de Derandomização) e a relação
$P\stackrel{?}{=}BPP\stackrel{?}{\subseteq}NP$ são intimamente
ligadas a geração pseudoaleatória e a $P\neq NP$.

\end{programs}

%% ---------------------------------------------------------------
\section{Linguagem FTM candidata}
%% ---------------------------------------------------------------

\begin{definition}[$\FTMDESCENT$]
Dada uma apresentação finita $\langle\Atlas\rangle$ de um atlas
federado trivalente, definir
\[
  \FTMDESCENT
  =\{\langle\Atlas\rangle : \exists\text{ dado global compatível
     com todos os cociclos}\}.
\]
\end{definition}

\begin{obligation}[C1 — Codificação]\label{ob:C1}
Fixar gramática, alfabeto e comprimento canônico de $\langle\Atlas\rangle$
de modo que a serialização seja verificável em $O(|\langle\Atlas\rangle|^k)$.
\end{obligation}

\begin{obligation}[C2 — Pertencimento a $NP$]\label{ob:C2}
Provar que o certificado (dado global) tem comprimento polinomial
e pode ser verificado em tempo polinomial.
Isso pode falhar com dados contínuos ou equivalências não efetivas.
\end{obligation}

\begin{obligation}[C3 — $NP$-dureza]\label{ob:C3}
Exibir redução many-one polinomial de um problema $NP$-completo
canônico para $\FTMDESCENT$, preservando respostas sim/não.
\end{obligation}

\begin{obligation}[C4 — Subclasses tratáveis]\label{ob:C4}
Identificar restrições (árvores, largura de treewidth limitada,
grupo finito fixo) que colocam o problema em $P$.
\end{obligation}

\begin{remark}
Mesmo que $\FTMDESCENT$ seja $NP$-completo (\OPEN),
isso \emph{não} separa $P$ de $NP$.
Uma separação requer argumento válido para \emph{todas} as máquinas
polinomiais e deve superar as barreiras da Seção~2A.
\end{remark}

%% ---------------------------------------------------------------
\section{Obrigações formais FTM para P vs NP}
%% ---------------------------------------------------------------

\begin{gap}[Barreira de relativização para FTM]
A construção federada usa oráculos implícitos (cociclos, descidas).
É necessário provar que o argumento não é relativizável no sentido
de Baker–Gill–Solovay. \OPEN
\end{gap}

\begin{gap}[Barreira de provas naturais para FTM]
Se a separação FTM usa propriedade combinatória de atlas,
deve-se verificar que tal propriedade não é ``natural''
no sentido de Razborov–Rudich. \OPEN
\end{gap}

\MATURITY{M0. A linguagem $\FTMDESCENT$ ainda carece de codificação
finita definitiva (C1 aberta). C2–C4 todas abertas.}

\end{document}
